%% APPENDIX D -- the false-alarm expectation, the control ensemble, and the
%% execution-block quality criterion.  Owner: r14-chance.  Carries
%% \label{app:falsealarm} (04_method.tex, 08_backmatter.tex) and
%% \label{sec:blockquality} (04_method.tex).
\section{The false-alarm expectation and the control ensemble}
\label{app:falsealarm}

This appendix gives the chance expectation for the unattributed crossings,
what it is conditioned on, and what the control ensemble calibrating it can
support.

\subsection{What the expectation is conditioned on}

\emph{One attribution set serves both sides of the comparison}, the set of
\S\ref{sec:technosearch}: the mask of Appendix~\ref{app:linemask} within
$\pm\MaskHalfKms$\,km\,s$^{-1}$ in the star's own rest frame, sulphur monoxide
excluded. So the \MkNSODecide{} sulphur monoxide coincidences are unattributed
crossings, and a control maximum on one of those transitions is an
unattributed exceedance. A systemic velocity is measured for the stars of
\CeNVsys{} of the census windows and the others are taken at rest, which moves
the masked share by \CeVsysSensPct{} points.

The expectation is built from the control positions, the only empirical null
the data provide, and each Class~A window contributes two quantities
measured on itself. The fraction of its own \NCtrl{} control positions that
reaches the trigger is the probability that a position in that field does so
-- no cell count enters it and no form is assumed for the tail -- and
summed over the windows \CeExpTrig{} are expected to trigger by chance. How
many of those would be attributed is measured on the channels the control
maxima land on, which the retained spectra record: \CeNExcAttr{} of
\CeNExc{} exceedances fall inside the mask, and the attribution therefore
takes \CeMaskPct{} per cent of the expectation. A maximum free to land
anywhere on its own grid would take \CeUnifPct{} per cent, the median window
offering \CeMaskPctMed{} per cent of its channels; the difference is that
\CeNFullMask{} narrow line-tuned windows lie wholly inside the mask and
their fields, which hold real emission, are where a control position most
often triggers, so a uniform landing would under-attribute them and leave
the unattributed expectation too high. What is
left is \CeExp{} expected unattributed Class~A crossings against \CeObs{}
observed.

Both sides run over the \CeNWin{} Class~A windows of the primary census, so
the expectation is a sum over rows and nothing is carried at an assumed rate.
The same comparison holds system by system, each system's expectation formed
from its own control positions: \MtNSysLowP{} of the \MtNSysChance{} systems
fall below one chance in twenty, against \MtNSysLowPExp{} expected. The system
holding the most crossings is \MtTopStar, toward which lie
\MtTopStarNCross{} of the \EvNCrossRaw{} crossings and \MtTopStarNUnattr{} of
the unattributed ones. Over its \MtTopSysNWin{} Class~A windows it gives
\MtTopSysObs{} against \MtTopSysExp{} expected. No system departs from its
own controls by more than chance allows.
A further \CeNWinTail{} Class~A windows arrived with the archival tail and lie
outside it, together with the \CeNTailCrossWord{} unattributed crossings
toward \CeTailStar{} they hold, which are reported beside the comparison and
enter neither side of it. No window holds more than one primary crossing, so
both sides count windows; the weaker cells behind those maxima are counted at
their own grain in \S\ref{sec:technosearch}.

The one assumption the searched data cannot test is that a star behaves like
a control position in its own field, and two samples measure it. In the census
the star carries \CeStep{} times as many above-threshold cells as a control
position. The ratio is measured between $T=\CeStepLevLo$ and
$T=\CeStepLevHi$, below the trigger, where no carrier claim is at stake, on
\CeStepNStar{} stellar cells over \CeStepNWin{} windows, and resampling the
execution blocks allows \CeStepLo--\CeStepHi{} at \CeCovPct{} per cent.
Read level by level the ratio rises from \CkRatioLo{} to \CkRatioHi{}
across that band. At the trigger itself it is \CeTrAt, but there it is formed
from the \CeStepNX{} above-threshold cells that are the crossings under test.
The step applied to the sum above is therefore the sub-trigger one, which
gives \CeExpX{} against the \CeObs{} observed. The reserved blocks measure
the same quantity on blocks the search did not reach. Their stars reach the
trigger in \StHoObs{} of \StHoNFine{} fine-channel windows, and the control
positions of those same windows predict \StHoExpFine, a ratio of
\StHoRatio. Every one of those
events lies in a reserved block that shares an observing session with a
searched one, and the \HsIndBlocks{} reserved blocks that share none hold
\HsIndObsWord{} trigger \HsIndObsNoun{} between them, so they neither confirm
the step nor contradict it (Appendix~\ref{sec:heldout}). In the \StHoNCoarse{} coarse-channel
reserved windows no star reaches the trigger, against \StHoExpCoarse{}
predicted. The same displacement is read in the body of the distribution,
where the stellar rank sits below one half (Appendix~\ref{app:radius}), and it
runs the same way.

A count is compared with the distribution of counts chance delivers, and
not with a range on its own mean. A rate is drawn from \CeNBoot{} resamplings
-- each window's rate redrawn from its own \NCtrl{} probes, the attributed
share from its own \CeNExc{} exceedances, and the execution blocks resampled
because windows share them and are not independent trials -- and the count
scatters about that rate. A count of \CeObs{} or more then arises with
probability \CeProbNull{} at the control positions alone, and \CeProbStep{}
with the step applied: the comparison is satisfied without the step rather
than by it. The largest added population that would leave the count still
consistent with chance at \CeCovPct{} per cent is \CeAddLim{} events, and
that is the limit on what the comparison can exclude. Read at the grain of
the individual cell rather than the window, the same accounting gives
\CkExp{} expected from the \CkNCtrlPerWin{} retained control probes of each
of those \CkNWin{} windows, \CkNCtrlPos{} positions in all, \CkExpSec{}
of them beside a window's largest. Against that it finds \CkNObs{} and
\CkNSecUnattr, with \CkLo--\CkHi{} covering \CeCovPct{} per cent of the same
resampling.

Every Class~B crossing in the survey comes from the one execution block of
\S\ref{sec:blockquality}, and over both classes together \CeExpSurv{} are
expected against the same \CeObs{} observed: a coarse single pointing holds
too few cells for its own field to reach the trigger by chance, and only
\NCtrlWinObs{} of the \NWindows{} windows hold a control position above
$T_{\rm trig}=\StageNullTrig$ at all. The Class~A comparison is the survey's.

\emph{The statistic is always the one measured over the window's whole field.}
\StNRep{} windows admit integrations inside the primary beam that a narrower
field would discard, recovering a median \StRepOnsrc{} times the on-source
time and moving the statistic by a median \StRepDT; the whole-field value is
the one adopted, on \StRepAdopted{} rows of Table~\ref{tab:ledger}, and each
window's control ensemble is measured with it, so neither side of the
comparison mixes two fields. Where the extraction catalogue's disposition
column disagrees -- on \StRepFell{} crossings, \StRepAttrClause{} -- the
ledger is the authority.

\subsection{Execution-block quality}
\label{sec:blockquality}

The comparison with the control positions assumes each block's ensemble
behaves like noise, and one block makes that false: every one of its \NCtrl{}
probes exceeds the trigger, the smallest at \BqFailCtrlMin. A criterion for
the condition is therefore applied to every block rather than to the one that
prompted it.

\emph{Two criteria are built from that statistic and they are not the same
test.} The statistic is the median of a window's own \NCtrl-position control
ensemble, which separates a field with a source in it from a field that is no
null at all. A \emph{window} is flagged when that median exceeds \BqFactor{}
the median of the same quantity over the survey, \DqThresh{} against
\DqSurvMed. The flag is a property of the window and cuts across attribution
by design: it fires on \DqNWin{} of the \NWindows{} windows in
\BkFlagNBlock{} blocks toward \DqNStar{} stars (\DqStars), and every one of
those windows carries a crossing, attributed and unattributed alike. A
\emph{block} is judged on the same median divided by the median over the
windows of its own class, since a maximum over a larger
channel$\times$drift grid sits higher and one threshold should mean the same
thing in both; a block fails when the median of that ratio over its windows
exceeds \BqFactor{} the scale of its class. The factor is the same in
both.

\emph{Neither criterion was fixed in advance.} The condition was noticed in
one block and the criterion written afterwards, so it is post hoc, and what
makes it trustworthy is how little the position of the line matters. The
failing block sits at \BqFailQ{} and the next most extreme at \BqNextQ, a
factor \QsRatio{} apart. Over \QsNStep{} thresholds spanning that interval
the same single block is selected every time, in number and in identity, so
the line may be drawn anywhere from \QsMarginLo{} to \QsMarginHi{} times
where it is. Outside that interval the answer does change, which is how one
can tell the rule discriminates: at \BqNextQ{} it selects \QsNAtNext{}
blocks, and at the median block, \QsQMed, \QsNAtMed{} of the
\BqNBlocks. Nor does the class treatment decide it. The same block is
selected with each class normalised on its own windows and with the survey
taken as one, and on this block's \QsFailNWinB{} coarse windows alone, which
are all the windows it has; and no block in the census is failed on one class
and passed on the other.

Applied to all \BqNBlocks{} blocks of the primary census, where the median
block sits at \BqQMed, the criterion fails \BqNFail: \BqFailBlock{} toward
\BqFailStar. Nothing else is close, and the unscaled form of the rule fails
the same \BqRawNFail{} block. Its \BqFailNWin{} windows are all
coarse-channel and
hold the survey's entire Class~B crossing population, so it carries
\EvNCrossFlag{} of the \EvNCrossRaw{} threshold crossings. \EvNRankFlag{}
of those outrank all \NCtrl{} control positions of their own window, and
no Class~A statistic here has ever contained it. Those \EvNCrossFlag{} appear
in Table~\ref{tab:ledger} with their block named and enter no count
elsewhere.

\subsection{The control ensemble: its calibration, and its resolution}

The ring geometry is set in \S\ref{sec:statistic}; its calibration and its
resolution are the subject here. Over the \VerNTested{} census windows that
retain a control vector the stellar ranks are consistent with the one half
exchangeability implies -- median \BkRankBalSurv{} with each star carrying
equal weight, \BkRankBalSurvLo{} to \BkRankBalSurvHi{} at \CeCovPct{} per
cent -- in both array
strata, the compact-array windows measured on each observation's own antenna
table rather than on a 12\,m beam. That is a weak statement and not a strong
one, since a shift smaller than about \VerKsPowerAca{} could not have been
seen in a sample of this size. The reserved blocks lean further from one half
when their windows are pooled, but those windows are concentrated on a few
stars: \BkRankTopFiveWin{} of \HoWin{} belong to five. The two
samples agree once each star carries equal weight, \BkRankBalHo{}
(\BkRankBalHoLo--\BkRankBalHoHi) against the \BkRankBalSurv{} above.
Compared star by star over the \BkRankPairedStars{} stars both samples hold
the difference is \BkRankPaired{} with \BkRankPairedLo{} to
\BkRankPairedHi{} allowed at the same coverage
(Appendix~\ref{sec:heldout}). What that settles is
the gap between the two samples' pooled medians, which is an effect of how
their windows are distributed over stars. It does not remove the displacement
itself: that is small, it is measured on the census's own cells at every level
up to the trigger, and the expectations above are stepped by it rather than
corrected for it.

\emph{The rank is not a significance test, and the reason is structural.}
Ranking the star against \NCtrl{} control positions of the same field
decides which windows deserve examination, but $1/(\NCtrl+1)=\RankFloorVal$
is the \emph{resolution} of the ensemble -- the smallest rank it can
express -- and not the probability that noise produced the window. A survey
of \NWindows{} windows would need a rank \BonfRankRatio{} times finer, and
correlation makes it coarser still: controls inside the synthesised beam are
not independent, so an annulus resolves no better than
\RankEffLo--\RankEffHi{} however many probes are drawn into it, with the
primary beam bounding how many a single pointing contains. That floor is the
design's central limitation: no single window here can be significant on the
rank alone, which is why confirmation requires independent recurrence.

\subsection{The trials budget}

The searched cells counted in \S\ref{sec:summary} are not independent
trials: channel smoothing and a drift grid stepped twice per channel
correlate neighbouring ones, which alone would leave \StGridFrac{} of them
independent. The control ensembles say otherwise, their maxima sitting at
\CtrlMaxRatioMed{} of the prediction for the maximum of the full searched
count treated as independent and Gaussian -- a budget predicting
\NCtrlWinExp{} windows with a control above $T_{\rm trig}$ against the
\NCtrlWinObs{} observed -- while fitting each window's effective count
against its own ensemble returns \StNindFrac{} of its searched cells. Nothing
turns on which is right, because the expectation above is summed from measured
rates and not from a cell count: by the cell-count route the number reaching
the trigger is \StExpModel{} rather than \CeExpTrig, \StModelAgreePct{} per
cent away.

\emph{The on-star crossings do exceed that budget, and the attribution
accounts for it.} A Gaussian tail over every searched cell predicts
\CellExpect{} stellar exceedances against the \NCross{} found. But \MkAttr{}
are attributed by the set above, \EvNCrossFlag{} come from the block of
\S\ref{sec:blockquality} and \CeNTailCross{} fall outside the census windows,
which leaves \CeObs{} against the \CeExpSurv{} predicted. No excess survives
the
attribution.

\emph{One stratum departs from uniformity, and the departure is where the
line emission is.} Count the windows whose stellar peak both reaches the
trigger and exceeds all \NCtrl{} control positions of its own window. The
$\NFine$ fine-channel windows hold \NStageOneFine{} against
\ExpStageOneFine{} expected; the coarse stratum, three times larger and
predicting \ExpStageOneCoarse, contains \NStageOneCoarse. A coarse
channel dilutes a narrow circumstellar line, so the class that can resolve
one is the class that finds it -- and an unresolved carrier would appear in
Class~A first for the same reason. Reaching the trigger is holding a
crossing, so each of those windows holds one and \CeFirstAttr{} hold a
crossing the mask attributes: the departure is a statement about the windows
carrying line emission and not about the fine stratum at large. Drop the
trigger condition and the star leads its own ring in \BkFirstBelow{} further
windows, at \BkFirstBelowTLo{} and \BkFirstBelowTHi. That is \NsrFirstObs{}
over all \NWindows{} windows against \NsrFirstExp{} expected. One quantity,
two conditions; the chance comparison uses the conditioned one.

\emph{What is left is bounded, and the bounds are null.} The real stellar
ranks agree at $\pv{\StarPseudoP}$ with the \NPseudo{} ranks each control
takes against the other probes of its own window; radial suppression is below
\RingRhoBoundPct{} per cent in the mean; and scrambling a window so nothing
coherent survives along a drift track carries \LNNWindows{} past the rank
floor and changes no disposition.

%% The chance expectation here is summed from the measured control exceedance
%% rates and each window's own masked grid share (chance_v414.py, round 310),
%% so no cell count and no Gaussian tail enters it and nothing in it is
%% conditioned on the crossings; the cell-count route is a cross-check only,
%% and the rank-conditioned expectation (\ClustStageMean against
%% \ClustStageNObs) is a pipeline diagnostic that must not be reachable as
%% "the" expectation.  tab:ctrlsep, tab:excess, tab:arraysplit and
%% fig:ctrldiag are generated and nothing cites them; r15-deposit, 2026-10-08:
%% they used to be described as "in the data release", and there is no data
%% release, so they are simply unpublished diagnostics.
%% tab:ledger lives here rather than in Sec. 5, beside the false-alarm
%% accounting that explains how many of its rows are expected.
%% The ledger is landscape because it is a wide one: fifteen columns, and
%% at the \scriptsize every other table in this paper is held to it measures
%% 605 pt against the 513 pt of a page's width and the 676 pt of its height.
%% Upright it fitted only at \tiny, which prints at 5.0 pt.
%% ★★★★ r18: ONE FLOAT, NOT TWO, AND THE COLUMN DEFINITIONS ARE NOW THIS
%% SUBSECTION RATHER THAN THE CAPTION.  A `sidewaystable*` costs a whole page
%% whatever its row count, so 56 rows in two parts cost two pages for content
%% that needs 507 pt of a page's 513.  The old caption ran to nineteen
%% rotated lines and -- because `rotating` centres a sideways double float in
%% \hsize, which this revtex-derived class leaves at the COLUMN width -- it
%% ran off the left edge of the paper with words missing from every line of
%% it.  The preamble patches the placement; `rotfit_r18.py` fails the build on
%% any ink outside the text block of a landscape page; and the method prose
%% that made the caption nineteen lines long is below, where a reader can
%% read it upright.
\subsection{Reading the crossing ledger}
\label{sec:ledgercolumns}

The columns of Table~\ref{tab:ledger} are the star, its execution block and
the ALMA band; the crossing frequency as the correlator delivered it, or,
where no peak frequency was measured, recovered from the line offset, and the
same frequency in the star's rest frame; the search statistic $T_\star$ and
the largest statistic $T_{\rm ctrl}$ among the window's \NCtrl{} control
positions; the nearest masked transition and the stellar-frame offset
$\Delta v_\star$ from it; and then the recurrence test -- $n_{\rm rep}$
covering repeat blocks spanning $\Delta t$ days, the statistic $T_{\rm pers}$
a carrier persisting at the discovery flux would have produced at the
predicted stellar-frame cell, the statistic $T_{\rm rep}$ measured there at
the carried drift, and the exclusion of Eq.~\ref{eq:excl}. Three of those
columns reproduce that exclusion, and two of them also give the limit defined
in Appendix~\ref{app:radius}.

Two kinds of blank appear in that test. Where the emission is present again in
a covering block it returns a recovery and not an exclusion, so
$T_{\rm rep}$ and $\sigma_{\rm excl}$ are blank and the disposition says which
it was. Rows whose window is data-quality flagged
(\S\ref{sec:blockquality}) are marked, and on those $T_{\rm pers}$ and the
exclusion are withheld instead: both are built from the discovery amplitude,
and in a window whose control ring stands as high as its star that amplitude
cannot be assigned to a point source at the star.

Not every interval in that column is measured from the blocks' own stamps.
Where the archive does not
index an execution block under its own identifier it carries the earliest
start time of its observing unit, so two blocks of one unit return one stamp
and the interval between them would be zero by construction. No such interval
is printed as a zero: a double dagger marks one measured from a unit's stamp
rather than from a block's own, and the blocks of
each unit so marked fall inside one day of each other on their own
barycentric terms, so an interval of tens of days is good to the day it is
printed to.

A crossing is attributed by the set defined above, and a dagger marks the
\MkNSODecide{} rows whose coinciding species that set excludes. The
transition named is the nearest in the star's frame, so where it is far from
every transition it is the magnitude of the offset and not the name that
carries the information. The delivered and rest-frame frequencies agree to
the digits printed in one row only, \MkCancStar, where the barycentric and
systemic terms -- \MkCancBaryMHz\,MHz each, in opposite senses -- cancel to
\MkCancNetkHz\,kHz. That star's systemic velocity,
\MkCancVSys\,km\,s$^{-1}$, is the barycentric correction at its epoch,
\MkCancVBary\,km\,s$^{-1}$, and its offset from \MkCancLine{} is
\MkCancDv\,km\,s$^{-1}$ either way. That star is also the fainter component
of a \CenPairSep{} arcsec pair. The
\CenPairNWin{} readings taken at its position are second extractions of the
same \CenPairNWin{} spectral windows of one execution block and not windows
of their own, so the per-window census counts them once, under
\CenPairNameA. Both components sit \CenPairOffA{} and \CenPairOffB{} arcsec
inside a \CenPairPbFwhmLo--\CenPairPbFwhmHi{} arcsec beam, with a response no
lower than \CenPairPbRespLo.

\begin{sidewaystable*}
\centering
\caption{\textbf{The crossing ledger, from which every disposition in this
paper can be reproduced.} All \EvNCrossRaw{} threshold crossings, one to a
window, including the \EvNCrossFlag{} in the failing block of
\S\ref{sec:bdblock}. The columns, the blanks and the dagger are defined
in \S\ref{sec:ledgercolumns}. Block identifiers omit the constant
\texttt{A002\_} prefix.}
\label{tab:ledger}
\scriptsize
%% GENERATED by ledger_v410.py -- do not hand-edit.
%% One row per threshold crossing, one epoch convention, stellar-frame attribution.
%% part 1 of 1, 56 rows; the split is made in the generator so the float cannot outgrow its page.
\begin{tabular}{@{}l@{\hspace{0.33em}}l@{\hspace{0.33em}}c@{\hspace{0.33em}}r@{\hspace{0.33em}}r@{\hspace{0.33em}}r@{\hspace{0.33em}}r@{\hspace{0.33em}}l@{\hspace{0.33em}}r@{\hspace{0.33em}}r@{\hspace{0.33em}}r@{\hspace{0.33em}}r@{\hspace{0.33em}}r@{\hspace{0.33em}}r@{\hspace{0.33em}}l@{}}
\hline
Star & block & B & $\nu_{\rm topo}$ (GHz) & $\nu_{\star}$ (GHz) & $T_\star$ & $T_{\rm ctrl}$ & line & $\Delta v_{\star}$ & $n_{\rm rep}$ & $\Delta t$ (d) & $T_{\rm pers}$ & $T_{\rm rep}$ & $\sigma_{\rm excl}$ & disposition \\
\hline
BD+05 1668 & Xc7fa6f\_X28a8 & 6 & 240.1016 & 240.1054 & 73.398 & 74.69 & CS($5$--$4$) & \ensuremath{-}5911.9 & 29 & 0.9--113 & n/a & +0.15 & n/a & unattributed; flagged \\
BD+05 1668 & Xc7fa6f\_X28a8 & 6 & 224.7422 & 224.7458 & 48.303 & 83.74 & H$_2$CO($3_{1,2}$--$2_{1,1}$) & \ensuremath{-}1264.5 & 29 & 0.9--113 & n/a & +0.37 & n/a & unattributed; flagged \\
BD+05 1668 & Xc7fa6f\_X28a8 & 6 & 226.7422 & 226.7458 & 41.355 & 75.48 & CN($2_{3/2}$--$1_{1/2}$) & +88.0 & 29 & 0.9--113 & n/a & +0.66 & n/a & unattributed; flagged \\
BD+05 1668 & Xc7fa6f\_X28a8 & 6 & 242.1953 & 242.1992 & 39.586 & 83.29 & CS($5$--$4$) & \ensuremath{-}3349.2 & 29 & 0.9--113 & n/a & +0.77 & n/a & unattributed; flagged \\
$\beta$~Pic & Xf5d76d\_Xeb8 & 3 & 115.2603 & 115.2711 & 30.765 & 7.34 & CO($1$--$0$) & \ensuremath{-}0.2 & 4 & 53$^{\ddagger}$ & n/a & -- & n/a & attributed; recurs \\
$\beta$~Pic & Xf5d76d\_Xe19 & 3 & 115.2603 & 115.2711 & 27.684 & 6.48 & CO($1$--$0$) & \ensuremath{-}0.2 & 4 & 53$^{\ddagger}$ & n/a & -- & n/a & attributed; recurs \\
HD 48370 & Xc26103\_X155a & 6 & 230.5117 & 230.5243 & 26.834 & 18.36 & CO($2$--$1$) & \ensuremath{-}17.8 & -- & -- & n/a & -- & n/a & attributed; flagged \\
$\beta$~Pic & Xf4df6f\_Xc7 & 3 & 115.2620 & 115.2711 & 17.907 & 5.76 & CO($1$--$0$) & \ensuremath{-}0.2 & 4 & 53 & n/a & -- & n/a & attributed; recurs \\
$\beta$~Pic & Xf5d76d\_Xcc5 & 3 & 115.2604 & 115.2711 & 17.288 & 5.52 & CO($1$--$0$) & \ensuremath{-}0.2 & 4 & 53$^{\ddagger}$ & n/a & -- & n/a & attributed; recurs \\
$\beta$~Pic & Xf5d76d\_X32f1 & 3 & 115.2606 & 115.2713 & 14.654 & 7.27 & CO($1$--$0$) & +0.4 & 4 & 53$^{\ddagger}$ & n/a & -- & n/a & attributed; recurs \\
$\beta$~Pic & Xd9668b\_X3a90 & 6 & 230.5161 & 230.5378 & 11.681 & 9.12 & CO($2$--$1$) & \ensuremath{-}0.3 & 5 & 1.0--1983 & n/a & -- & n/a & attributed; recurs \\
$\beta$~Pic & X6f1341\_X1484 & 6 & 230.5284 & 230.5379 & 9.832 & 8.20 & CO($2$--$1$) & \ensuremath{-}0.1 & 5 & 25--1984 & n/a & -- & n/a & attributed; recurs \\
$\beta$~Pic & X7116f1\_X1785 & 6 & 230.5272 & 230.5384 & 9.677 & 14.13 & CO($2$--$1$) & +0.5 & 5 & 25--1959 & n/a & -- & n/a & attributed; recurs \\
$\beta$~Pic & Xd9668b\_Xa9df & 6 & 230.5160 & 230.5377 & 8.948 & 8.41 & CO($2$--$1$) & \ensuremath{-}0.4 & 5 & 1.0--1984 & n/a & -- & n/a & attributed; recurs \\
$\beta$~Pic & Xa7a216\_X23e4 & 6 & 230.5273 & 230.5377 & 7.988 & 6.03 & CO($2$--$1$) & \ensuremath{-}0.4 & 5 & 601--1313 & n/a & -- & n/a & attributed; recurs \\
61 Vir & Xc079b5\_X82f & 7 & 346.3138 & 346.3232 & 5.958 & 5.65 & SO($9_{8}$--$8_{7}$) & \ensuremath{-}177.6 & 3 & 0.08--3.0 & 9.7 & +0.39 & 4.9 & unattributed; no recurrence \\
HD 285968 & Xd98580\_X3bf6 & 6 & 230.5021 & 230.5447 & 5.956 & 8.12 & CO($2$--$1$) & +8.7 & 3 & 13--44 & n/a & -- & n/a & attributed; recurs \\
HD 285968 & Xdb7ab7\_X5b4e & 6 & 230.5116 & 230.5447 & 5.868 & 7.22 & CO($2$--$1$) & +8.7 & 3 & 5.0--44 & n/a & -- & n/a & attributed; recurs \\
HD 285968 & Xdb4b9a\_X108f & 6 & 230.5100 & 230.5447 & 5.835 & 7.94 & CO($2$--$1$) & +8.7 & 3 & 5.0--39 & n/a & -- & n/a & attributed; recurs \\
CP$-$72 2713 & Xff0235\_X4a6d & 7 & 344.2697 & 344.2965 & 5.809 & 5.68 & SO($8_{8}$--$7_{7}$) & \ensuremath{-}12.3$^{\dagger}$ & 5 & 0.09--2.1 & 6.3 & \ensuremath{-}0.82 & 4.8 & unattributed; no recurrence \\
HD~139084~B & Xcd8029\_Xb6b0 & 7 & 345.8241 & 345.8241 & 5.763 & 17.70 & CO($3$--$2$) & +24.4 & -- & -- & -- & -- & -- & attributed \\
HD 14055 & Xfff3d8\_Xd46 & 7 & 331.7699 & 331.7664 & 5.575 & 5.99 & $^{13}$CO($3$--$2$) & +1068.7 & 30 & 0.9--650 & 170.4 & \ensuremath{-}0.02 & 5.6 & unattributed; no recurrence \\
HD 23484 & X10b22d7\_X2445 & 6 & 230.7904 & 230.8006 & 5.520 & 5.55 & CO($2$--$1$) & +341.5 & 4 & 0.10--2.1 & 10.4 & +1.25 & 4.3 & unattributed; no recurrence \\
HD 161868 & Xff45a6\_X10ca8 & 7 & 345.6338 & 345.6491 & 5.481 & 6.16 & SO($2_{3}$--$2_{1}$) & \ensuremath{-}48.0$^{\dagger}$ & 13 & 0.02--6.1 & 92.2 & +0.24 & 5.5 & unattributed; no recurrence \\
HIP 10679 & Xa95c04\_X14cd & 6 & 220.4185 & 220.4035 & 5.456 & 6.76 & $^{13}$CO($2$--$1$) & +6.5 & -- & -- & -- & -- & -- & attributed \\
HD~14055 & X11adad7\_X25c5e & 7 & 345.4564 & 345.4303 & 5.448 & 6.82 & SO($2_{3}$--$2_{1}$) & \ensuremath{-}237.8 & 30 & 0.05--653 & 14.0 & +0.03 & 5.1 & unattributed; no recurrence \\
HD 207129 & Xc19d6f\_X5706 & 6 & 230.4226 & 230.4070 & 5.433 & 5.52 & CO($2$--$1$) & \ensuremath{-}170.3 & 4 & 0.8--425 & 69.7 & \ensuremath{-}0.93 & 5.5 & unattributed; no recurrence \\
$\beta$~Pic & Xd9668b\_X3a90 & 6 & 241.7512 & 241.7739 & 5.420 & 5.81 & CS($5$--$4$) & \ensuremath{-}3869.7 & 1 & 1.0 & 5.7 & +0.72 & 3.4 & unattributed; no recurrence \\
HR~1010 & Xc6fa08\_X1053 & 6 & 230.5133 & 230.5287 & 5.410 & 5.65 & CO($2$--$1$) & \ensuremath{-}12.0 & 7 & 2.0--14 & 8.4 & $<$+4.60 & $>$2.1 & attributed; not recovered \\
ALMA J153702653-33192492 & Xff0235\_X8392 & 6 & 230.5222 & 230.5392 & 5.405 & 13.41 & CO($2$--$1$) & +1.6 & 6 & 18--141 & 6.2 & $<$+4.15 & $>$1.4 & attributed; not recovered \\
HD 110411 & Xe45e29\_Xb7d4 & 6 & 219.4538 & 219.4357 & 5.393 & 5.85 & C$^{18}$O($2$--$1$) & \ensuremath{-}170.2 & 1 & 10 & 6.7 & +0.62 & 3.8 & unattributed; no recurrence \\
HD 31392 & X129754b\_X6946 & 6 & 231.1998 & 231.2231 & 5.289 & 5.81 & H($30\alpha$) & \ensuremath{-}876.3 & 8 & 0.06--41 & 15.7 & \ensuremath{-}0.26 & 5.1 & unattributed; no recurrence \\
HD 14055 & Xff99e1\_X24be0 & 7 & 331.7078 & 331.7035 & 5.277 & 5.69 & $^{13}$CO($3$--$2$) & +1011.6 & 30 & 0.9--652 & 252.2 & +0.15 & 5.3 & unattributed; no recurrence \\
AU Mic & Xa42f75\_X319 & 6 & 230.6864 & 230.6705 & 5.274 & 6.07 & CO($2$--$1$) & +172.3 & 2 & 310--455 & 6.4 & \ensuremath{-}0.16 & 4.2 & unattributed; no recurrence \\
HD 69830 & X12209f7\_X16edb & 7 & 322.4325 & 322.4444 & 5.259 & 5.42 & C$^{18}$O($3$--$2$) & \ensuremath{-}6268.5 & 5 & 0.07--29 & 8.9 & \ensuremath{-}0.48 & 4.8 & unattributed; no recurrence \\
ALMA J153702653-33192492 & Xfe62c1\_X398 & 6 & 216.1844 & 216.2036 & 5.228 & 6.02 & SiO($5$--$4$) & \ensuremath{-}1244.6 & 7 & 18--126 & 11.0 & +1.03 & 4.3 & unattributed; no recurrence \\
AU Mic & X899169\_X1838 & 6 & 230.8252 & 230.8294 & 5.221 & 5.85 & CO($2$--$1$) & +379.0 & 2 & 145--310 & 10.0 & +0.19 & 4.5 & unattributed; no recurrence \\
HD 14055 & Xff99e1\_X2c3bb & 7 & 345.5256 & 345.5216 & 5.186 & 5.66 & SO($2_{3}$--$2_{1}$) & \ensuremath{-}158.6 & 30 & 0.09--651 & 172.8 & +2.00 & 5.1 & unattributed; no recurrence \\
HD 14055 & Xfffde1\_X3835 & 7 & 345.0099 & 345.0070 & 5.177 & 6.56 & SO($2_{3}$--$2_{1}$) & \ensuremath{-}604.8 & 30 & 0.07--649 & 187.2 & \ensuremath{-}1.78 & 5.2 & unattributed; no recurrence \\
HD 14055 & Xff99e1\_X2b2ee & 7 & 344.7007 & 344.6965 & 5.175 & 5.92 & SO($8_{8}$--$7_{7}$) & +336.0 & 30 & 0.09--651 & 210.7 & \ensuremath{-}0.00 & 5.2 & unattributed; no recurrence \\
LHS 1140 & Xd9d7c7\_X6fd2 & 6 & 230.6355 & 230.6268 & 5.136 & 5.68 & CO($2$--$1$) & +115.5 & 4 & 5.0--27 & 7.3 & +0.33 & 4.0 & unattributed; no recurrence \\
HD 10647 & Xff294f\_Xaf70 & 7 & 331.6570 & 331.6943 & 5.120 & 5.91 & $^{13}$CO($3$--$2$) & +1003.3 & 10 & 0.1--592 & 125.8 & +0.39 & 5.1 & unattributed; no recurrence \\
CD$-$57 1054 & Xcf525f\_X5aa4 & 7 & 346.5099 & 346.5265 & 5.108 & 5.71 & SO($9_{8}$--$8_{7}$) & \ensuremath{-}1.7$^{\dagger}$ & 1 & 39 & 6.0 & \ensuremath{-}0.27 & 4.1 & unattributed; no recurrence \\
ALMA J153702653-33192492 & Xf8f6a9\_X10359 & 6 & 217.7292 & 217.7269 & 5.091 & 5.41 & H$_2$CO($3_{0,3}$--$2_{0,2}$) & \ensuremath{-}680.5 & 7 & 3.8--242 & 11.6 & +0.07 & 4.6 & unattributed; no recurrence \\
GJ 849 & Xd9668b\_X841a & 6 & 230.8769 & 230.8581 & 5.090 & 5.83 & CO($2$--$1$) & +416.3 & 5 & 0.05--51 & 9.3 & \ensuremath{-}1.69 & 5.3 & unattributed; no recurrence \\
ALMA J153702653-33192492 & Xfe62c1\_X398 & 6 & 218.8729 & 218.8923 & 5.090 & 5.56 & H$_2$CO($3_{2,1}$--$2_{2,0}$) & +181.2 & 7 & 18--126 & 10.5 & \ensuremath{-}1.04 & 5.0 & unattributed; no recurrence \\
HD 14055 & X1001fdc\_X29ae & 7 & 331.1559 & 331.1546 & 5.082 & 6.26 & $^{13}$CO($3$--$2$) & +513.8 & -- & -- & -- & -- & -- & unattributed; no spectrum \\
$\beta$~Pic & Xd9668b\_Xa9df & 6 & 241.8367 & 241.8594 & 5.079 & 5.99 & CS($5$--$4$) & \ensuremath{-}3765.1 & 1 & 1.0 & 4.8 & \ensuremath{-}0.63 & 3.9 & unattributed; no recurrence \\
ALMA J153702653-33192492 & Xf8f6a9\_X10359 & 6 & 218.6711 & 218.6688 & 5.048 & 5.77 & H$_2$CO($3_{2,1}$--$2_{2,0}$) & \ensuremath{-}125.1 & 7 & 3.8--242 & 11.6 & +0.00 & 4.6 & unattributed; no recurrence \\
HD 69830 & X12277ed\_X178c6 & 7 & 321.2005 & 321.2149 & 5.046 & 6.13 & C$^{18}$O($3$--$2$) & \ensuremath{-}7387.7 & 5 & 6.0--36 & 10.2 & \ensuremath{-}0.46 & 4.7 & unattributed; no recurrence \\
HR 1010 & Xc66ce1\_Xaa9 & 6 & 230.8348 & 230.8491 & 5.044 & 5.72 & CO($2$--$1$) & +404.6 & 7 & 1.1--14 & 10.3 & \ensuremath{-}1.53 & 5.2 & unattributed; no recurrence \\
61 Vir & Xc067f7\_X822 & 7 & 345.5599 & 345.5678 & 5.038 & 5.79 & SO($2_{3}$--$2_{1}$) & \ensuremath{-}118.5 & 3 & 2.0--3.0 & 9.3 & +0.45 & 4.2 & unattributed; no recurrence \\
HD 170773 & Xff0235\_X8c42 & 7 & 331.0101 & 331.0234 & 5.029 & 5.75 & $^{13}$CO($3$--$2$) & +394.8 & 4 & 0.9--113 & 74.8 & +1.56 & 4.9 & unattributed; no recurrence \\
ALMA J153702653-33192492 & Xd9d7c7\_Xbb64 & 7 & 347.2179 & 347.1890 & 5.009 & 5.68 & SiO($8$--$7$) & \ensuremath{-}122.2 & 1 & 0.05 & 5.0 & \ensuremath{-}2.06 & 5.0 & unattributed; no recurrence \\
HD~14055 & X1190de8\_Xd526 & 7 & 344.4676 & 344.4515 & 5.006 & 5.96 & SO($8_{8}$--$7_{7}$) & +122.7 & 30 & 20--611 & 11.6 & \ensuremath{-}0.60 & 4.8 & unattributed; no recurrence \\
$\eta$~Crv & X122b6ff\_X1041e & 7 & 357.2815 & 357.2475 & 5.006 & 5.50 & HCO$^{+}$($4$--$3$) & +431.3 & 1 & 87 & 11.0 & \ensuremath{-}0.62 & 4.8 & unattributed; no recurrence \\
\hline
\end{tabular}

\end{sidewaystable*}

%% R2-5: the positive controls of the recurrence stage, with the recovering block,
%% its session, its epoch, the separation and the statistic there.  Generated by
%% confirm_r19.py (round 806); the main text cites it once, in Sec. 5.6.
\begin{table*}
\centering
\caption{\textbf{The recurrence stage's positive controls, and the baseline
each was recovered over.} The \QcNCtl{} line-attributed crossings the
criterion of \S\ref{sec:summary} recovers, toward \QcNStar{} stars: for each,
the block whose own crossing satisfies it at the same stellar-frame
frequency, that block's epoch, its separation from the discovery block and
the statistic it reaches there. Block identifiers omit the constant
\texttt{A002\_} prefix and are printed as session and block, so the first
field is the observing session: every recovering block here belongs to a
different session from its own discovery block and lies more than \EpSep{}
day from it, which is what the criterion demands. $T$ is the statistic of the
recovering block's own crossing at that frequency and not a single-cell
reading of a repeat, which is the quantity the \RcNExcl{} exclusions of
Table~\ref{tab:ledger} are built from. $\ddagger$~marks a separation measured
from an observing unit's stamp rather than from a block's own.
\QcNUnrecWord{} of these crossings are present in blocks of their own
observing unit as well, whose internal separations the archive does not record;
Table~\ref{tab:ledger} carries the same mark on those.}
\label{tab:confirm}
\small
%% GENERATED by confirm_r19.py -- do not hand-edit.
%% The recurrence stage's positive controls and the baseline each was recovered over.
\begin{tabular}{@{}llrllrrr@{}}
\hline
Star & transition & $\nu_\star$ (GHz) & discovery block & recovering block & epoch & $\Delta t$ (d) & $T$ \\
\hline
$\beta$~Pic & CO(1--0) & 115.2711 & \texttt{Xf5d76d\_Xcc5} & \texttt{Xf4df6f\_Xc7} & 2022-01-09 & 53$^\ddagger$ & 17.91 \\
$\beta$~Pic & CO(1--0) & 115.2711 & \texttt{Xf5d76d\_Xe19} & \texttt{Xf4df6f\_Xc7} & 2022-01-09 & 53$^\ddagger$ & 17.91 \\
$\beta$~Pic & CO(1--0) & 115.2711 & \texttt{Xf5d76d\_Xeb8} & \texttt{Xf4df6f\_Xc7} & 2022-01-09 & 53$^\ddagger$ & 17.91 \\
$\beta$~Pic & CO(1--0) & 115.2711 & \texttt{Xf4df6f\_Xc7} & \texttt{Xf5d76d\_Xeb8} & 2022-03-02 & 53$^\ddagger$ & 30.76 \\
$\beta$~Pic & CO(1--0) & 115.2713 & \texttt{Xf5d76d\_X32f1} & \texttt{Xf4df6f\_Xc7} & 2022-01-09 & 53 & 17.91 \\
$\beta$~Pic & CO(2--1) & 230.5377 & \texttt{Xa7a216\_X23e4} & \texttt{Xd9668b\_X3a90} & 2019-03-12 & 1312 & 11.68 \\
$\beta$~Pic & CO(2--1) & 230.5377 & \texttt{Xd9668b\_Xa9df} & \texttt{X6f1341\_X1484} & 2013-10-06 & 1984 & 9.83 \\
$\beta$~Pic & CO(2--1) & 230.5378 & \texttt{Xd9668b\_X3a90} & \texttt{X6f1341\_X1484} & 2013-10-06 & 1983 & 9.83 \\
$\beta$~Pic & CO(2--1) & 230.5379 & \texttt{X6f1341\_X1484} & \texttt{Xd9668b\_X3a90} & 2019-03-12 & 1983 & 11.68 \\
$\beta$~Pic & CO(2--1) & 230.5384 & \texttt{X7116f1\_X1785} & \texttt{Xd9668b\_X3a90} & 2019-03-12 & 1958 & 11.68 \\
HD~285968 & CO(2--1) & 230.5447 & \texttt{Xdb7ab7\_X5b4e} & \texttt{Xd98580\_X3bf6} & 2019-03-15 & 44 & 5.96 \\
HD~285968 & CO(2--1) & 230.5447 & \texttt{Xdb4b9a\_X108f} & \texttt{Xd98580\_X3bf6} & 2019-03-15 & 39 & 5.96 \\
HD~285968 & CO(2--1) & 230.5447 & \texttt{Xd98580\_X3bf6} & \texttt{Xdb7ab7\_X5b4e} & 2019-04-28 & 44 & 5.87 \\
\hline
\end{tabular}

\end{table*}
